\section{Discussion}\label{sec:discussion}

\subsection{Interpretation}
In our runs, the slice did not act as one large device. The model used one chip unless the program mapped work to the
others. This behavior is expected from how JAX places work on devices. Our measurements quantify the effect for
AlphaFold2 on a TPU v5e slice. On a rented slice of eight chips, a program that uses one chip leaves seven chips idle
(Section~\ref{sec:res-cost}).
At 118 residues and up to eight queries per call, batching on one chip with \texttt{vmap} did not replace multi-query \texttt{pmap}. % src: recomputed: INPUT.toy_length

A user with many independent queries can run one query on each chip in parallel, as multi-query \texttt{pmap} did. Each
\texttt{pmap} call ran queries of one length, so batches of mixed lengths need padding or grouping by length, as
ParaFold does~\cite{zhong2022parafold}. We did not compare \texttt{pmap} with eight separate processes that each use one
chip. Many users run AlphaFold2 on several accelerators in this way. Our \texttt{pmap} runs used JAX's implementation of
\texttt{pmap} with \texttt{shard\_map}, but we did not call \texttt{jax.shard\_map} directly or shard a batch with
\texttt{jax.jit}.

Automatic partitioning did not shorten the call for one query. Our program gave the partitioner no sharding
annotations, and we found none in the source of AlphaFold2. We did not try to add annotations or to change the model.
The external ensemble ran eight featurizations of one query on eight chips. The external ensemble averaged one output
with a collective operation. With random weights, this output was zero. We did not test whether this average helps
prediction.

\subsection{Limitations}\label{sec:limitations}
The model ran with random weights. The trained weights have the same shapes, so we expect the same steady-state times.
The parameters are an argument of the compiled program, so XLA does not specialize the program on their values. We did
not measure the times with the trained weights. Every MSA held only the query, and no run relaxed the predicted
structure with a force field. We did not measure deep MSAs or the models with templates. Only the recycling sweep ran with recycles. We
timed the calls of the model alone, so our costs cover only part of the work for one query.
We measured one TPU generation, the three models without templates, and inputs of at most 1000 residues. % src: results/sweep/scaling_law_data.json -> grid[*].length

Most settings ran once. Hoefler and Belli ask authors to report confidence intervals for nondeterministic
measurements~\cite{hoefler2015benchmarking}. Papadopoulos et al. make repeated experiments the first of their principles
for performance evaluation in the cloud~\cite{papadopoulos2021principles}. Most of our settings meet neither of these
recommendations. We measured the TPU in one period only.
In the three repeated runs, the coefficient of variation was at most 1.00\% for every timed call. % src: recomputed: RR.steady_state_seconds.cv_pct, RR.init_params_seconds.cv_pct, RR.first_predict_seconds.cv_pct
This check does not cover changes between Jobs or between sessions.

Times also differ between experiments at similar lengths.
At 120 residues, the repeated-alphabet sequence took 0.408~s in \texttt{predict}. % src: results/sweep/sequence_length_sweep.json -> [1].num_residues, [1].steady_state_seconds
At 118 residues, the fixed sequence took 0.469~s to 0.471~s in \texttt{predict}. % src: recomputed: INPUT.toy_length, TPU.single_query_steady_min_max
With \texttt{vmap} and $B$ = 1, the repeated-alphabet sequence of 118 residues took 0.4875~s in \texttt{apply}, although \texttt{predict} includes \texttt{apply}. % src: results/sweep/batching_sweep.json -> [0].batch_size, [0].steady_state_seconds; recomputed: INPUT.toy_length
These runs differ in Job, script, or input family. We did not study these differences. Between experiments, we draw no
conclusion from discrepancies of this size. On the T4 and on the CPU runtime, the changes between sessions were large.

The records have gaps.
The August Colab runs used a copy of the script, and this copy exists only inside the notebooks of commit 2755250. % src: recomputed: V2.first_commit
We cannot identify the script of the baseline run on the TPU. The TPU logs are not in the repository. The sweep files are
summaries of the Job output. The committed TPU Jobs pinned only JAX.
